\chapter{Phát hiện trang đã được truy cập: \texttt{pgaccess()}}

\newpage

\section{Ý tưởng}

RISC-V ghi nhận việc truy cập ở cấp trang bằng bit \texttt{PTE\_A} (bit 6): mỗi khi bộ dịch địa chỉ nạp một PTE vào TLB, phần cứng bật bit này. Vì vậy chỉ cần đọc và \textbf{xoá} \texttt{PTE\_A} là biết được, trong khoảng thời gian giữa hai lần gọi, những trang nào đã được đọc hoặc ghi. Đây là cơ chế mà các bộ thu gom rác dùng để tìm trang ``còn sống''.

System call có ba tham số, đúng như đề yêu cầu:

\begin{center}
\texttt{int pgaccess(void *base, int len, void *mask)}\quad (\texttt{user/user.h:28})
\end{center}

\noindent trong đó \texttt{base} là địa chỉ ảo của trang đầu, \texttt{len} là số trang cần kiểm tra, và \texttt{mask} là địa chỉ vùng nhớ người dùng nhận bitmask: bit $i$ ứng với trang \texttt{base + i*PGSIZE}, bit thấp nhất là trang đầu tiên.

\section{Đường đi của system call}

\begin{table}[H]
    \centering
    \renewcommand{\arraystretch}{1.3}
    \setlength{\tabcolsep}{8pt}
    \begin{tabular}{|l|l|l|}
        \hline
        \rowcolor[rgb]{0.8, 0.9, 1.0}
        \textbf{Tầng} & \textbf{Tệp} & \textbf{Nội dung} \\ \hline
        Số hiệu & \texttt{kernel/syscall.h:25} & \texttt{SYS\_pgaccess = 24} \\ \hline
        Bảng dispatch & \texttt{kernel/syscall.c:106,134} & prototype và \texttt{sys\_pgaccess} \\ \hline
        Bảng tên & \texttt{kernel/syscall.c:161} & \texttt{"pgaccess"} cho \texttt{trace} \\ \hline
        Cài đặt & \texttt{kernel/sysproc.c:199-256} & \texttt{sys\_pgaccess()} \\ \hline
        Stub user & \texttt{user/usys.pl:41} & \texttt{entry("pgaccess")} \\ \hline
    \end{tabular}
    \caption{Các điểm sửa để thêm system call \texttt{pgaccess}}
\end{table}

Bảng tên được cập nhật vì đồ án 1 đã thêm system call \texttt{trace}; nếu thiếu, gọi \texttt{trace(1)} sẽ in tên rỗng thay vì \texttt{pgaccess}.

\section{Cài đặt}

\begin{codebox}[sysPgaccess]{\texttt{sys\_pgaccess()}}
    \begin{minted}{c}
// The bitmask is a single uint64, so one call can report at most 64 pages.
#define PGACCESS_MAX 64

// Report which of the len pages starting at base have the RISC-V accessed
// bit set, then clear that bit.  Bit i of the mask stands for page
// base + i*PGSIZE.
//
// Returns -1 if the range is invalid or only partly mapped.  The whole
// range is checked and the mask is copied out before any access bit is
// cleared, so a rejected call leaves the page table untouched.
uint64
sys_pgaccess(void)
{
  struct proc *p = myproc();
  uint64 base, maskaddr, bits = 0;
  pte_t *pte;
  int len;

  argaddr(0, &base);
  argint(1, &len);
  argaddr(2, &maskaddr);

  if(len <= 0 || len > PGACCESS_MAX)
    return -1;
  if(base >= MAXVA || base % PGSIZE != 0 ||
     base + (uint64)len * PGSIZE > MAXVA)
    return -1;
  if(maskaddr >= MAXVA || maskaddr + sizeof(bits) > MAXVA)
    return -1;

  // Pass 1: every page must be a valid user mapping.
  for(int i = 0; i < len; i++){
    if((pte = walk(p->pagetable, base + i * PGSIZE, 0)) == 0)
      return -1;
    if((*pte & PTE_V) == 0 || (*pte & PTE_U) == 0)
      return -1;
  }

  // Pass 2: collect the access bits.  Do not clear them yet: if copyout()
  // fails, the next call must still see the same access state.
  for(int i = 0; i < len; i++){
    pte = walk(p->pagetable, base + i * PGSIZE, 0);
    if(*pte & PTE_A)
      bits |= 1UL << i;
  }

  if(copyout(p->pagetable, maskaddr, (char *)&bits, sizeof(bits)) < 0)
    return -1;

  // Pass 3: clear only the bits that were successfully reported.
  for(int i = 0; i < len; i++){
    if(bits & (1UL << i)){
      pte = walk(p->pagetable, base + i * PGSIZE, 0);
      *pte &= ~PTE_A;
    }
  }
  return 0;
}
    \end{minted}
\end{codebox}

Năm quyết định thiết kế:

\begin{enumerate}[itemsep=2.5pt, topsep=2.5pt, leftmargin=1.5cm]
    \item \textbf{Bitmask nằm trong kernel.} \texttt{bits} là biến cục bộ trên kernel stack; kernel điền xong mới \texttt{copyout()} một lần. Nếu ghi trực tiếp vào con trỏ người dùng thì mỗi bit là một lần \texttt{copyout()} và không thể kiểm tra con trỏ trước khi đã ghi dở.
    \item \textbf{Giới hạn 64 trang.} Một \texttt{uint64} chỉ có 64 bit, nên \texttt{len > 64} bị từ chối thay vì âm thầm cắt bớt. Đề cho phép đặt giới hạn trên; nếu cần nhiều hơn thì phải truyền mảng \texttt{uint64}.
    \item \textbf{Kiểm tra địa chỉ trước khi đi bộ bảng trang.} \texttt{base} phải căn trang, toàn bộ dải phải nằm dưới \texttt{MAXVA}, và vùng \texttt{mask} cũng phải nằm trong không gian địa chỉ hợp lệ. Nhờ vậy không có lời gọi nào chạm tới \texttt{walk()} với địa chỉ không hợp lệ.
    \item \textbf{Kiểm tra trước, chỉ xoá sau khi \texttt{copyout()} thành công.} Pha 1 kiểm tra toàn bộ dải; pha 2 chỉ đọc các bit \texttt{PTE\_A}; pha 3 mới xoá những bit đã được báo. Nếu dải có trang chưa ánh xạ hoặc con trỏ \texttt{mask} không ghi được, lời gọi trả $-1$ và \textbf{không} làm mất trạng thái truy cập của các trang.
    \item \textbf{Kiểm tra biên trước khi \texttt{walk()}.} \texttt{walk()} \texttt{panic} khi \texttt{va >= MAXVA}. Vì vậy biểu thức \texttt{base + len*PGSIZE} được kiểm tra trước, tránh cả trường hợp tràn số $2^{64}$.
\end{enumerate}

Cờ \texttt{PTE\_A} được thêm vào \texttt{kernel/riscv.h:365}:

\begin{codebox}[pteA]{Cờ \texttt{PTE\_A}}
    \begin{minted}{c}
#define PTE_A (1L << 6) // accessed
    \end{minted}
\end{codebox}

\section{Chương trình kiểm thử}

\texttt{user/pgtbltest.c} gồm \textbf{28 check}, chia thành năm nhóm:

\begin{table}[H]
    \centering
    \renewcommand{\arraystretch}{1.3}
    \setlength{\tabcolsep}{8pt}
    \begin{tabular}{|l|c|p{8.5cm}|}
        \hline
        \rowcolor[rgb]{0.8, 0.9, 1.0}
        \textbf{Nhóm} & \textbf{Số check} & \textbf{Nội dung} \\ \hline
        Trạng thái ban đầu & 3 & Cấp 32 trang bằng \texttt{sbrk()}; lần gọi đầu trả 0 vì kernel vừa xoá trang qua direct map\\ \hline
        Ghi và đọc & 7 & Ghi trang 0, đọc trang 5: mask đúng hai bit; gọi lần ba trả 0 vì bit đã bị xoá; trang cuối dải cũng được báo \\ \hline
        Dải ngắn & 2 & \texttt{len = 4} chỉ xét 4 trang đầu, bỏ qua các trang đã truy cập phía sau \\ \hline
        Đầu vào sai & 12 & Dọn bit còn lại; con trỏ mask xấu và \texttt{base} lệch trang bị từ chối mà không xoá bit; dải đã \texttt{sbrk} âm (chưa ánh xạ), \texttt{len = 0}, \texttt{len} âm, \texttt{len > 64}, \texttt{base = MAXVA} --- tất cả trả $-1$ \\ \hline
        In theo yêu cầu & 4 & \texttt{vmprint()} trả 0, heap vẫn dùng được, bảng trang vẫn \texttt{walk} được, và trang mới vẫn ghi nhận bit \texttt{A} sau khi in \\ \hline
    \end{tabular}
    \caption{Các nhóm kiểm thử trong \texttt{user/pgtbltest.c}}
\end{table}

Ba điểm đáng chú ý trong cách kiểm thử:

\begin{itemize}[itemsep=0pt, topsep=2.5pt, leftmargin=1.5cm]
    \item Test \textbf{không} giả định trạng thái ban đầu của \texttt{PTE\_A}. Lời gọi đầu tiên vừa kiểm tra vừa xoá mọi bit còn sót, nên kết quả các lần sau là tất định dù trước đó có truy cập gì.
    \item Các test từ chối kiểm tra thêm rằng lời gọi lỗi không xoá bit: sau khi cố tình truy cập một trang, gọi \texttt{pgaccess()} với con trỏ mask xấu hoặc \texttt{base} lệch trang phải trả $-1$, rồi lời gọi hợp lệ ngay sau đó vẫn thấy bit \texttt{A} của trang đó.
    \item Trường hợp ``dải chưa ánh xạ'' không dùng một địa chỉ cố định nào, mà \texttt{sbrk(-32*PGSIZE)} để thu hồi đúng 32 trang vừa cấp. Nhờ vậy phép thử không phụ thuộc vào bố cục địa chỉ cụ thể của xv6.
\end{itemize}

\section{Kết quả và hạn chế}

\begin{codebox}[pgtbltest]{\texttt{pgtbltest} trên QEMU}
    \begin{minted}{text}
pgtbltest: page-table tests
PASS: pgaccess: heap grows by NPAGES pages
PASS: pgaccess: first call succeeds
PASS: pgaccess: freshly mapped pages are not reported
PASS: pgaccess: untouched page reads as zero
PASS: pgaccess: second call succeeds
PASS: pgaccess: exactly the written and read pages are reported
PASS: pgaccess: third call succeeds
PASS: pgaccess: the previous call cleared the access bits
PASS: pgaccess: last page call succeeds
PASS: pgaccess: last page of the range is reported
PASS: pgaccess: short range call succeeds
PASS: pgaccess: short range ignores pages past len
PASS: pgaccess: range cleared
PASS: pgaccess: bad mask pointer is rejected
PASS: pgaccess: call after bad mask succeeds
PASS: pgaccess: bad mask does not clear access bits
PASS: pgaccess: unaligned base is rejected
PASS: pgaccess: call after unaligned base succeeds
PASS: pgaccess: unaligned base does not clear access bits
PASS: pgaccess: unmapped range is rejected
PASS: pgaccess: len 0 is rejected
PASS: pgaccess: negative len is rejected
PASS: pgaccess: more than MAXPAGES pages is rejected
PASS: pgaccess: virtual address beyond MAXVA is rejected
PASS: vmprint: dumps the page table on demand
PASS: vmprint: the heap is still usable afterwards
PASS: vmprint: the page table is still walkable
PASS: vmprint: the new page still records accesses
ALL TESTS PASSED
    \end{minted}
\end{codebox}

Bốn dòng cuối thuộc \texttt{vmprint\_test()} (mục 4.4); phần bảng trang mà \texttt{vmprint()} in ra được lược bớt ở đây cho gọn.

Ngoài ra \texttt{usertests -q} trên cùng bản build cũng cho \texttt{ALL TESTS PASSED}, nên việc thêm số hiệu syscall 24--25 và đụng vào PTE không làm hỏng các bài kiểm thử có sẵn của xv6.

Ba hạn chế đã biết:

\begin{itemize}[itemsep=0pt, topsep=2.5pt, leftmargin=1.5cm]
    \item \textbf{Giới hạn 64 trang và một \texttt{uint64}.} Muốn quét nhiều hơn phải đổi ABI thành mảng bitmask, khi đó \texttt{len} và kích thước vùng \texttt{mask} phải khớp nhau.
    \item \textbf{Chỉ xoá \texttt{PTE\_A}, không đụng \texttt{PTE\_D}} (dirty, bit 7). \texttt{pgaccess} trả lời ``đã truy cập'', không phân biệt đọc với ghi; muốn phân biệt phải đọc thêm \texttt{PTE\_D} và định nghĩa lại ABI.
    \item \textbf{Phụ thuộc phần cứng đặt \texttt{PTE\_A} tự động.} QEMU bật bit này trong lúc đi bộ bảng trang. Trên CPU chỉ hỗ trợ \texttt{Svade} (fault khi \texttt{A = 0}), xv6 không có handler nên sẽ panic; bài lab giả định môi trường QEMU.
\end{itemize}
